% 第五章 實驗結果
\chapter{實驗結果}\label{ch:results}

本章分三部分。第一部分（第 \ref{sec:mech}、\ref{sec:gemm} 節）說明省電的機制與效益的上界；第二部分（第 \ref{sec:hpl}–\ref{sec:gp} 節）是具正確性檢查的完整程式：HPL、HPL-MxP、Quantum ESPRESSO、DF-MP2 與高斯過程地統計；第三部分（第 \ref{sec:crossover}–\ref{sec:ai} 節）是尾數歸零與改變格式做法（INT8 模擬、單精度、FP8）的分界。除非特別說明，百分比都是與同一功耗上限、同一節點上的原生版本比較。

% ==================================================================
\section{省電機制：位元不翻轉或數值為零}\label{sec:mech}

\subsection{填值對照}\label{sec:mech-fill}

尾數歸零同時改變兩件事：低位元不再隨資料翻轉，而且被捨去部分的數值變成 0。為了分開這兩種效應，我們把被捨去的位元改填全 1，或填入一組所有元素共用的固定非零樣式；這兩種填法同樣不會翻轉，但數值不是 0。我們在 1000~W 下對三個層次各量測一次：FP64 tensor core 矩陣乘法、只使用一般 FP64 運算單元的乘加迴圈，以及純記憶體搬移（圖~\ref{fig:mech}）。填值版本都以最高時脈 1980~MHz 執行，因此比較的是功耗；原生 FP64 矩陣乘法仍有 95\% 的時間撞上限、時脈只有 1822~MHz，其功耗是被上限截住的值，所以「不翻轉」所佔的比例會算得偏保守。

\begin{figure}[tbp]
\centering
\figfile{fig5_01_mechanism}
\caption[省電機制：被捨去的位元填 0、全 1 或固定樣式時的省電]{省電機制：保留 0 位、被捨去的位元填 0、全 1 或固定樣式時，三個層次在 1000~W 下比原生少用的模組功耗。填 1 與固定樣式同樣不翻轉，但數值不為 0；括號內為佔填 0 省電的比例。}\label{fig:mech}
\begin{figsource}Slurm 工作 1752–1753、1762、1815、1858；\path{gh200/results/m1_toggle_summary.csv} 等\end{figsource}
\end{figure}

在這組量測中（1000~W 下每個組合各量一次），三個層次的方向一致：只要讓低位元不再翻轉，就取得完整省電的 62–97\%；填 0 比填 1 或固定樣式再多省一些。記憶體搬移幾乎完全是「不翻轉」的效應（86–97\%）；乘法器則還有額外的「數值為 0」效應：保留 20 位時，填 1 的低位與另一個運算元相乘的部分積仍會隨後者翻轉，填 0 則整片不動。保留 0 位時兩個運算元的尾數都是常數，部分積本來就不翻轉，但填 0 仍比填 1 省電（707~W 對 774~W），可能來自乘積的對齊與累加較單純（全 0 時乘積是 2 的冪），此點尚未驗證。在 500~W（功耗受限、比較效能；單一節點、各兩次）重做同樣的對照，「不翻轉」佔效益的 77–92\%，方向相同。不過樣本很少（1000~W 各一次、500~W 一個節點兩次，填值只做了保留 0 位與 20 位），這些比例應視為這組量測的估計，而非精確的分解。此結果在 GH200 的模組功耗上為前人提出的翻轉假說~\cite{Gregersen2024}提供了量化的支持；資料搬移的部分也與 DRAM 能耗隨翻轉率變化的觀察一致~\cite{Chatterjee2017,Ghose2018}。

\subsection{不只是 tensor core：非矩陣乘法的 FP64 程式}\label{sec:nongemm}

為了確認省電機制不限於 tensor core，我們也測試不使用矩陣乘法的 FP64 程式；這也是 Ozaki 類模擬無法取代的部分，因為模擬只適用於稠密矩陣乘法。我們先在核心層級測試：只用一般 FP64 運算單元的乘加迴圈、記憶體複製與 triad、N 體重力計算，以及一個 CG 稀疏求解器（圖~\ref{fig:nongemm}）。

\begin{figure}[tbp]
\centering
\figfile{fig5_09_nongemm}
\caption[非矩陣乘法核心在 1000~W 下的省電]{非矩陣乘法核心在 1000~W 下保留 36 位與 20 位時，比原生減少的模組功耗。}\label{fig:nongemm}
\begin{figsource}Slurm 工作 1815、1858、1860；\path{gh200/results/m4_*_summary.txt}\end{figsource}
\end{figure}

在低功耗上限下，這些省電會轉成加速：乘加迴圈在 500~W 快 4.8–9.0\%（36–20 位），記憶體複製在 450~W 快 7.7–14.0\%。但 N 體計算只快 1–4\%：它第一步就是座標相減 $\mathit{dx} = x_j - x_i$，差值的尾數又是滿的，後面約 20 個運算照常翻轉。這說明效益隨「計算深度」遞減，只有直接讀取截斷資料的運算受惠。CG 幾乎沒有效益，因為被截斷的矩陣值只佔每圈記憶體流量的一部分，索引與向量不能截斷。

\paragraph{第三方真實軟體}為了確認核心層級的結果能否延伸到真實程式，我們在三個不以矩陣乘法為主的 FP64 GPU 程式中，於設定階段對「儲存後每一步都會被乘進去的係數陣列」歸零一次（不設定時與原版逐位元相同），結果列於表~\ref{tab:apps}。

\begin{table}[tbp]
\centering
\caption{第三方非矩陣乘法程式的係數截斷結果}\label{tab:apps}
\small
\begin{tabularx}{\textwidth}{@{}>{\raggedright\arraybackslash}p{3.0cm}>{\raggedright\arraybackslash}p{2.6cm}>{\centering\arraybackslash}p{1.9cm}L@{}}
\toprule
軟體 & 被截斷的係數 & 係數佔記憶體流量 & 結果 \\
\midrule
MFEM（免組裝有限元素，3827 萬未知數） & 每個求積點的幾何與係數資料 & 約 43\% & 450~W 快 0.9\%（36 位）／2.0\%（20 位），550~W 快 0.7\%／1.4\%；解的相對變化 $\le \ee{2.3}{-11}$ \\
SW4（地震波，1.26 億格點） & 材料參數 $\rho$、$\mu$、$\lambda$ & 約 9\% & 350~W 快 0.6\%（36 位）／1.3\%（20 位），450~W 快 0.3\%／0.5\%；波形最大相對差 $\le \ee{3.4}{-6}$ \\
LAMMPS（EAM 分子動力學） & 位能表 & 約 0\%（常駐快取） & 各上限 0.0–0.4\%，落在雜訊內（負面對照成立） \\
\bottomrule
\end{tabularx}
\end{table}

MFEM 的係數佔了近一半的記憶體流量、直接進入乘法器，收益仍只有 1–2\%。三者的收益順序（MFEM $\ge$ SW4 $>$ LAMMPS）與係數佔記憶體流量的比例一致。

\paragraph{小結}尾數歸零的省電不限於 tensor core，一般 FP64 運算單元與記憶體搬移也能省電；但在真實程式中，被截斷的資料只是耗電的一部分，核心層級的 7–18\% 會被其餘滿尾數的運算稀釋到 1–2\%。

% ==================================================================
\section{功耗上限下的矩陣乘法微基準}\label{sec:gemm}

我們先在最單純的工作負載上量出效益的上界。測試使用 cuBLAS 矩陣乘法（FP64 為 $32768\times32768\times2048$，BF16 為 $32768\times32768\times6144$），輸入是 $[-1,1)$ 的均勻隨機值；在 3 個節點、300–1000~W 共八個上限下，每個上限以交錯順序各執行 2 次（圖~\ref{fig:gemm}）。

\begin{figure}[tbp]
\centering
\figfile{fig5_02_gemm_caps}
\caption[矩陣乘法微基準在各功耗上限下的加速]{FP64 與 BF16 矩陣乘法在各功耗上限下保留不同位數的加速（95\% 信賴區間，3 個節點各 2 對；FP64 保留 27 位與 BF16 的點分別向左、右錯開約 7~W 以免重疊）。}\label{fig:gemm}
\begin{figsource}Slurm 工作 1667–1669（work1–3）；\path{gh200/results/e1_ci.csv}\end{figsource}
\end{figure}

大多數信賴區間只有 $\pm$0.1–1.2 個百分點，效益遠大於量測雜訊。換算成等效功耗，FP64 保留 27 位在 500–800~W 相當於多給 79–162~W，BF16 保留 0 位相當於多給 121~W 到 200~W 以上。最佳能效也隨之提高：FP64 從原生的 82~GFLOP/J（500~W）升到保留 27 位的 90~GFLOP/J（400~W），BF16 從 950~GFLOP/J（600~W）升到保留 0 位的 1159~GFLOP/J（500~W）。效益與功耗上限並非單調關係，在 300–400~W 最大。一個可能的解釋（尚未驗證）是：上限極低時電壓已在下限，功耗大致與時脈成正比，省下的功耗能換到較多時脈；而在 1000~W 時，FP64 與 BF16 矩陣乘法仍受功耗上限限制，所以仍有 3.9\%（FP64 保留 35 位）到 15.7\%（BF16 保留 0 位）的加速。

需要強調的是，這些是均勻隨機資料、保留位數很低時的數字，屬於翻轉最劇烈的情況，代表效益的上界，而不是應用的操作點。下一節檢視具正確性檢查的完整程式。

% ==================================================================
\section{HPL 與 HPL-MxP：具正確性檢查的完整基準}\label{sec:hpl}

\subsection{實驗設定與原生 HPL 的調校}\label{sec:hpl-tune}

HPL 的實驗使用其內建的量化選項：在每次 FP64 矩陣乘法之前，以一個捨入核心把兩個輸入捨入到 $k$ 位（平手取偶的最近捨入），乘法本身仍是 FP64 cuBLAS。實驗使用 8 個節點、$N = 286720$（調參後的設定），在 400–1000~W 下，每個上限在同一個工作內依「原生、36 位、28 位、28 位、36 位、原生」的交錯順序執行，即每個變體兩次、與原生配成兩對。

為了確保比較的基準是最佳的原生效能，我們先調校原生 HPL。最早的量測沿用 H200 時期的設定（$N = 262144$、$\mathit{NB} = 1024$、$2\times4$ 網格、$f = 0.05$、NCCL 預設值），1000~W 下原生只有 374.8–376.0~TFLOP/s。我們在 1000~W 掃描 $N$、$\mathit{NB}$、$P\times Q$、$f$ 與 NCCL 參數（共 24 組設定、26 次執行）：NCCL 固定 8 個通道與 2~MiB 緩衝貢獻約 4.5\%，$N$ 加大到 286720 再多約 2\%，$4\times2$ 網格反而慢 14\%。同一個最佳設定重跑三次為 390.2–398.8~TFLOP/s（相差約 2\%；其中 390.2 與 393.4 在同一個工作內，398.8 在另一個工作），$f$、NCCL 通道與緩衝、$\mathit{NB}$、$N$ 的其他鄰近設定都沒有超出這個範圍（$f = 0.15$ 反而慢 3\%）。以下的功耗上限掃描都使用這個設定，在同一組 8 個節點上進行（各結果對應的 Slurm 工作列於附錄 \ref{app:jobs}）。

\subsection{功耗上限下的效能}\label{sec:hpl-perf}

表~\ref{tab:hpl-perf} 與圖~\ref{fig:hpl} 為調參後的結果。在規格表可設範圍（450–1000~W）內的 500–800~W，保留 36 位有效位數（含隱藏位，即 35 個尾數位元）快 4.0–5.0\%，每次求解的模組耗電少 4.1–5.1\%；換算成等效功耗，500–700~W 相當於每個節點多給 42–68~W。800~W 落在 800 與 1000~W 兩個量測點之間，只能直線內插；效能對上限的曲線向下彎，內插值（127~W，若以實測功耗為橫軸則為 103~W）必然偏高，所以不列入範圍。在低於規格表下限的 400~W，加速達 11.0\%、耗電少 9.3\%，但這一點的 SM 時脈只多 6.3\%，而 500–800~W 時脈多 4.9–6.2\%、效能多 4.0–5.0\%，兩者相當；400~W 的額外收益來自時脈以外的因素，原因尚未查明，因此我們把它與 500–800~W 分開報告。

等效功耗與百分比的關係取決於原生曲線的斜率：上限低時曲線陡，多給 1~W 換到的效能多，所以 400~W 的 11.0\% 只相當於 39~W；上限高時曲線變平，同樣幾個百分點的效能要多給更多瓦數才補得回來。因此不同上限的 \DW{} 不能只看百分比比較。

量測雜訊方面，調參時同一設定的 390.2–398.8~TFLOP/s（約 2\%）主要是不同工作之間的差異，而且調參掃描是依序執行各設定、沒有配對。上限掃描則在同一個工作內以交錯順序配對：在 400–800~W，同一變體兩次之間只差 0.03–0.6\%，兩對各自的加速也一致（表~\ref{tab:hpl-perf} 的「兩對」欄，相差最多 0.6 個百分點），比 4–5\% 的效應小一個數量級。不過每格只有兩對、都在同一組節點上，我們不給信賴區間，也不能排除這組節點本身的特性；增加輪數與節點組是第 \ref{sec:limits} 節所列的待補工作。

在預設的 1000~W 上限下，調參後的原生 HPL 仍有 75\% 的取樣撞上限（SM 平均 1936~MHz），截斷後只剩 7\%（SM 1977~MHz），平均快 3.9\%，模組功耗從 930~W 降到 886~W，每次求解的模組耗電少 7.5\%（驅動程式讀數）。但 1000~W 的兩次原生相差 2.2\%，兩對的加速為 +2.7\% 與 +5.1\%、省電為 6.3\% 與 8.7\%，所以 1000~W 的數字不確定性較大。未調參的設定（$N = 262144$）在 400–800~W 快 3.4–8.3\%、1000~W 快 0.7\% 並省電 7.7\%，方向相同。

\begin{table}[tbp]
\centering
\caption[調參後 8 節點 HPL 在各功耗上限下的效能]{調參後 8 節點 HPL 在各功耗上限下的效能（$N = 286720$，TFLOP/s，兩次的平均）}\label{tab:hpl-perf}
\small
\begin{tabularx}{\textwidth}{@{}lCCCCCC@{}}
\toprule
每節點模組上限 & 原生 & 保留 36 位 & 加速 & 兩對的加速（\%） & 等效多出的上限 \DW & 每次求解的耗電 \\
\midrule
400~W$^{*}$ & 225.0 & 249.9 & +11.0\% & 11.1、11.0 & 39~W & $-9.3\%$ \\
500~W & 288.0 & 300.7 & +4.4\% & 4.4、4.4 & 42~W & $-4.1\%$ \\
600~W & 318.0 & 334.0 & +5.0\% & 5.0、5.1 & 57~W & $-5.1\%$ \\
700~W & 346.2 & 362.9 & +4.8\% & 4.5、5.1 & 68~W & $-4.8\%$ \\
800~W & 370.8 & 385.6 & +4.0\% & 4.1、3.9 & （127~W）$^{\dagger}$ & $-4.2\%$ \\
1000~W & 393.9 & 409.2 & +3.9\% & 2.7、5.1 & — & $-7.5\%$ \\
\bottomrule
\end{tabularx}
\begin{tablenote}「36 位」為含隱藏位的有效位數，即 35 個尾數位元。$^{*}$低於規格表的可設範圍（450–1000~W）。$^{\dagger}$在 800 與 1000~W 之間直線內插，偏高，不列入範圍。耗電為驅動程式讀數的模組耗電。資料來源：\path{gh200/results/e3_tuned_summary.csv}、\path{gh200/results/e3_tuned_raw.csv}。截斷後的數字為研究用途，並非正式成績。\end{tablenote}
\end{table}

\begin{figure}[tbp]
\centering
\figfile{fig5_03_hpl}
\caption[調參後 8 節點 HPL 的效能與功耗上限]{調參後 8 節點 HPL 的效能與功耗上限：原生、保留 36 位與保留 28 位。保留 28 位的殘差已接近門檻，只作為失效邊界的參考。}\label{fig:hpl}
\begin{figsource}Slurm 工作 3315（調參後的功耗上限掃描）、1820（多組 $N$）；\path{gh200/results/e3_tuned_summary.csv}\end{figsource}
\end{figure}

依規格，每顆 GH200 的 FP64 tensor core 峰值為 67~TFLOP/s~\cite{NVIDIAGH200}，8 顆合計約 536~TFLOP/s。調參後的原生 HPL 在 1000~W 達到約 73\%（重測兩次平均 393.9~TFLOP/s，調參時單次最高 398.8~TFLOP/s），保留 36 位達到 409.2~TFLOP/s（研究用途，並非正式成績）。保留 28 位可以再多 0–5 個百分點（例如 700~W 時 372.3~TFLOP/s），但它在這個 $N$ 的殘差已達 15.4（門檻 16），在 $N = 196608$ 時更是失敗，所以不作為操作點。

\subsection{殘差與失效邊界}\label{sec:hpl-resid}

HPL 最後會檢查縮放殘差是否小於 16。殘差會隨矩陣大小 $N$ 變化，所以我們在 1000~W 以三組 $N$ 檢查（每組 2 次，殘差完全重現），再加上調參後的 $N = 286720$（功耗上限掃描的 12 次，殘差也完全重現），結果列於表~\ref{tab:hpl-resid}。

\begin{table}[tbp]
\centering
\caption{HPL 縮放殘差隨保留位數與矩陣大小 $N$ 的變化（門檻 16）}\label{tab:hpl-resid}
\small
\begin{tabularx}{\textwidth}{@{}lccccL@{}}
\toprule
保留位數 & $N = 131072$ & $N = 196608$ & $N = 262144$ & $N = 286720$ & 判定 \\
\midrule
原生（53 位） & \ee{1.5}{-4} & \ee{1.3}{-4} & \ee{2.6}{-5} & \ee{7.6}{-5} & 通過 \\
36 & 0.16 & 0.072 & 0.042 & 0.038 & 通過，距門檻約百倍（最小 98 倍） \\
28 & 13.97 & \textbf{18.09} & 14.70 & 15.40 & $N = 196608$ 失敗，其餘都接近門檻 \\
24 & 456 & 462 & 91 & — & 全部失敗 \\
\bottomrule
\end{tabularx}
\end{table}

從 36 位減到 28 位少了 8 位，殘差放大 86–410 倍，量級與 $2^{8} = 256$ 相符。這表示 28 位距門檻不到 1 位，換一個 $N$ 就可能失敗；36 位則有約 7 位的餘裕，因此我們以 36 位作為 HPL 的操作點。

\subsection{HPL-MxP}\label{sec:mxp}

HPL-MxP 以 BF16 完成 LU 分解，再以 FP64 迭代修正。我們在每次寫回 BF16 時歸零，使用 8 個節點、$N = 491520$。此處使用的測試矩陣是對角優勢矩陣（對角元等於該列非對角元絕對值的和），對角 tile 使用近似分解（\texttt{--approx-diag}），迭代修正使用 Richardson 法搭配 FP32 預條件子。依 HPL-MxP 的規則，嚴格對角優勢的矩陣只能用於測試~\cite{HPLMxPRules}；在強對角優勢的矩陣上，很粗的 $L$、$U$ 仍是好的預條件子，所以以下「保留 3 位仍穩定」的結論不能直接外推到正式允許的矩陣。

\begin{figure}[tbp]
\centering
\figfile{fig5_04_hplmxp}
\caption[8 節點 HPL-MxP 在各功耗上限下的加速]{8 節點 HPL-MxP 保留 3 位與 0 位在各功耗上限下的加速（對角優勢的測試矩陣，規則只允許用於測試）。}\label{fig:mxp}
\begin{figsource}Slurm 工作 1671；\path{gh200/results/e2_summary.csv}\end{figsource}
\end{figure}

如圖~\ref{fig:mxp} 與表~\ref{tab:mxp-perf} 所示，保留 3 位在 $N = 245760$–491520 的四組大小都穩定：殘差 1.4–3.8，固定 3 次修正。在 $N = 491520$，它在 500~W 快 6.5\%，600–800~W 只快 1.8–2.5\%；600–800~W 的兩對加速都落在 1.5–2.8\%、方向一致，但兩對不足以給出有意義的信賴區間，我們不宣稱 600~W 以上的差異顯著。在低於規格表下限的 400~W 快 12.7\%。等效功耗在 400–700~W 為 34–45~W，耗電量少 3.4–10.7\%。每個上限的修正次數都從原生的 1 次變成 3 次；LU 本身的加速從 400~W 的 16.0\% 降到 800~W 的 6.0\%，扣除多出的修正時間後，600~W 以上只剩約 2\%。保留 0 位雖然快 7.1–26.0\%，但 $N = 491520$ 的殘差已達 10.6–11.1（門檻 16），$N = 393216$ 時為 11.7–13.5，在 $N = 245760$ 時修正次數暴增到 52 次、效能下降 70\%，所以只是失效邊界。

\begin{table}[tbp]
\centering
\caption[8 節點 HPL-MxP 在各功耗上限下的效能]{8 節點 HPL-MxP 在各功耗上限下的效能（$N = 491520$，TFLOP/s，報告值含迭代修正，兩次的平均）}\label{tab:mxp-perf}
\small
\begin{tabularx}{\textwidth}{@{}lCCCCCC@{}}
\toprule
每節點模組上限 & 原生 BF16 & 保留 3 位 & 加速 & 兩對的加速（\%） & 等效多出的上限 \DW & 每次求解的耗電 \\
\midrule
400~W$^{*}$ & 2076 & 2340 & +12.7\% & 11.8、13.7 & 35~W & $-10.7\%$ \\
500~W & 2823 & 3007 & +6.5\% & 6.8、6.2 & 45~W & $-6.6\%$ \\
600~W & 3229 & 3310 & +2.5\% & 2.2、2.8 & 37~W & $-3.7\%$ \\
700~W & 3449 & 3519 & +2.0\% & 1.7、2.4 & 34~W & $-4.7\%$ \\
800~W & 3655 & 3720 & +1.8\% & 2.1、1.5 & — & $-3.4\%$ \\
\bottomrule
\end{tabularx}
\begin{tablenote}$^{*}$低於規格表的可設範圍（450–1000~W）。800~W 的等效功耗需要 800~W 以上的原生點，而原生在 1000~W 會觸發機箱的硬體功耗煞車，所以不報；1000~W 的比較也不報。資料來源：\path{gh200/results/e2_summary.csv}、\path{gh200/results/e2_raw.csv}。\end{tablenote}
\end{table}

保留 0 位在這個 $N$ 上更快（例如 400~W 時 2616~TFLOP/s），但在較小的 $N$ 上修正次數暴增，所以不作為操作點。另有一個附帶發現：風扇全開時，原生 HPL-MxP 在 1000~W 會觸發機箱電源的硬體功耗煞車（每次 4–10 秒），反而比 800~W 慢；尾數歸零版本的尖峰功耗較低，在開啟計數器的診斷工作中，保留 0 位的版本沒有觸發。這個現象只在一個診斷工作中觀察到，因此只作為附帶觀察。

\subsection{小結}

在具殘差檢查把關的完整基準上，功耗受限時尾數歸零都帶來正面收益；但保留位數必須以多組問題大小驗證，並選擇有餘裕的操作點；HPL-MxP 的結果目前限於對角優勢的測試矩陣。

% ==================================================================
\section{科學軟體：Quantum ESPRESSO 與 PySCF}\label{sec:sci}

\subsection{所需的尾數位數}\label{sec:sci-bits}

我們以攔截函式庫截斷所有 FP64 矩陣乘法，觀察物理結果（總能量）在多少位時開始偏離，結果列於表~\ref{tab:bits}。

\begin{table}[tbp]
\centering
\caption{科學軟體所需的尾數位數（攔截所有 FP64 矩陣乘法）}\label{tab:bits}
\small
\begin{tabularx}{\textwidth}{@{}>{\raggedright\arraybackslash}p{5.2cm}L@{}}
\toprule
系統 & 結果 \\
\midrule
Quantum ESPRESSO 金表面（112 原子） & 最近捨入 24 位誤差 \ee{8.1}{-6}~Ry、20 位 \ee{5.5}{-5}~Ry、16 位發散 \\
Quantum ESPRESSO 矽（Si 1000） & 28 位誤差 \ee{3.5}{-7}~Ry；20 位是失效邊緣（一次收斂、一次中止） \\
PySCF 纈胺黴素 DFT（門檻 $10^{-7}$） & 32–40 位都與原生同樣 12 圈收斂；32 位誤差 \ee{1.2}{-5}\Eh（約 0.007~kcal/mol） \\
PySCF 紫杉醇 HF+MP2 & 往零截斷 44 位誤差 \ee{5.6}{-8}\Eh；36 位 \ee{2.7}{-5}\Eh \\
\bottomrule
\end{tabularx}
\end{table}

平面波密度泛函只需要 20–28 位，高斯基底的量子化學需要 32–36 位；後者與 Dawson 等人以模擬得到的「約 37 位有效尾數即已足夠」一致~\cite{Dawson2024}。

\begin{figure}[tbp]
\centering
\figfile{fig5_05_qe_rounding}
\caption[Quantum ESPRESSO 金表面的總能量誤差與保留位數]{Quantum ESPRESSO 金表面的總能量誤差與保留位數：往零截斷與最近捨入。}\label{fig:qe}
\begin{figsource}Slurm 工作 1729、1738（金表面攔截掃描）\end{figsource}
\end{figure}

\subsection{捨入方式的影響}\label{sec:sci-rounding}

如圖~\ref{fig:qe} 所示，同樣保留 32 位與 28 位時，往零截斷的總能量誤差分別是最近捨入的 35 倍與 81 倍，遠大於單次矩陣乘法的差距（式~\eqref{eq:errmodel} 的係數 $\alpha$ 只差約 2.6 倍）：截斷永遠把數值往 0 拉，誤差有固定方向，會在自洽迭代中一路累積；最近捨入沒有偏差，誤差會互相抵銷。

\subsection{效能與耗電}\label{sec:sci-perf}

這些程式在 1000~W 下用不到上限，所以我們在會碰到上限的 450–550~W 量測，結果列於表~\ref{tab:qe-perf}。每個系統有兩個工作，各以交錯順序執行；每個工作的第一次執行（原生）是 GPU 的暖機，偏慢（例如金表面 96.9 秒，其餘原生為 92.2–94.7 秒），會使含它的那一對收益偏大。表中排除每個工作的第一對、以其餘 4 對計算；這條規則是看過資料後才訂的，但不依賴哪一對比較慢。用全部 6 對的結果列於表註，差異在 0.3 個百分點以內，結論相同。更乾淨的做法是在交錯順序之前先執行一次不計入的暖機，這是待補的實驗。

\begin{table}[tbp]
\centering
\caption{Quantum ESPRESSO 程式內截斷與自動調整位元數的時間與耗電}\label{tab:qe-perf}
\small
\begin{tabularx}{\textwidth}{@{}L>{\raggedright\arraybackslash}p{1.3cm}>{\centering\arraybackslash}p{1.25cm}>{\centering\arraybackslash}p{2.2cm}>{\centering\arraybackslash}p{2.2cm}>{\raggedright\arraybackslash}p{2.6cm}@{}}
\toprule
做法 & 系統 & 上限 & 時間 & 耗電量 & 總能量誤差 \\
\midrule
依收斂門檻自動調整位元數 & 金表面 & 450~W & $-4.0 \pm 0.4\%$ & $-3.7 \pm 0.5\%$ & 0（印出的每一位都相同） \\
依收斂門檻自動調整位元數 & 矽 & 550~W & $-2.4 \pm 0.2\%$ & $-2.4 \pm 0.2\%$ & 0（印出的每一位都相同） \\
程式內固定 24 位 & 金表面 & 450~W & $-2.4 \pm 0.5\%$ & $-2.4 \pm 0.4\%$ & \ee{3}{-8}~Ry \\
程式內固定 28 位 & 矽 & 550~W & $-2.2 \pm 0.3\%$ & $-2.2 \pm 0.5\%$ & $<10^{-8}$~Ry \\
攔截所有 GEMM，最近捨入 28 位（函式庫內建時的上限） & 矽 & 550~W & $-6.2\%$ & $-6.4\%$ & \ee{3.5}{-7}~Ry \\
\bottomrule
\end{tabularx}
\begin{tablenote}前四列為程式內的做法，與未修改的原生程式交錯執行、成對比較；數字為排除每個工作含暖機的第一對後，4 對的平均 $\pm$ 95\% $t$ 區間。用全部 6 對時依序為 $-4.2 \pm 0.7$／$-3.9 \pm 1.1$、$-2.4 \pm 0.5$／$-2.5 \pm 0.6$、$-2.7 \pm 1.1$／$-2.7 \pm 1.6$、$-2.3 \pm 0.6$／$-2.3 \pm 0.6\%$（時間／耗電）。最後一列與同樣付出複製成本的基準比較，各執行一次。耗電為驅動程式讀數的模組耗電。資料來源：\path{gh200/results/qe_adapt_au_6pairs.csv}、\path{gh200/results/qe_adapt_si_6pairs.csv}。\end{tablenote}
\end{table}

金表面上，自動調整在 6 對中每一對都比固定 24 位快 0.6–1.8 個百分點；矽上也每對都略快於固定 28 位，但只差 $0.18 \pm 0.14$ 個百分點。因此自動調整至少追平事後才找得到的最佳固定位數，而且不需要事先試驗。最後一列攔截所有矩陣乘法，代表「如果函式庫內建截斷」時的上限；攔截本身約有 5\% 的複製成本。

自動調整的規則（式~\eqref{eq:adapt}）讓位元數隨收斂從 16 升到 39 位，使截斷誤差大致不超過 Davidson 當圈本來就容許的誤差，並由後續迭代修正；因此總能量在印出的所有位數（$10^{-8}$~Ry）都與原生相同，SCF 圈數也相同。不過金表面的收益有一部分並非功耗效應：自動調整讓 Davidson 少做了約 7–9\% 的工作（\texttt{h\_psi} 呼叫 311 次對 336 次）。在一次功耗上限被外部重設為 1000~W、不受功耗限制、時脈相同的執行中，它仍比原生與固定 24 位快約 2\%（73.5–73.7 秒對 74.8–75.1 秒），與 450~W 時領先固定 24 位的幅度相同。所以金表面的 4.0\% 中約有一半來自迭代減少，原因尚未查明；矽上沒有這個現象（206 次對 205 次），2.4\% 可視為功耗效應。

\subsection{單精度的可行性}\label{sec:sci-fp32}

Quantum ESPRESSO 的運算元值域放得進 FP32，所以我們也嘗試把所有 GEMM 改為 FP32、TF32 或 3xTF32（3xTF32 是攔截工具內自行撰寫的實作：拆成兩個 TF32、以三次 FP32 累加的 tensor core 乘法計算）。結果金表面的 FP32 在第 26 圈、TF32 與 3xTF32 在第 2 圈中止；矽的 FP32 在第 14 圈、3xTF32 在第 9 圈中止，積分電荷也明顯偏離。原因在第 \ref{sec:crossover-fp32} 節說明。這是把所有 GEMM 都替換掉的極端做法；只在部分步驟使用 FP32 的混合精度（例如 Woo 等人~\cite{Woo2023}）不在此限。

\subsection{小結}

科學軟體需要 20–36 位，在功耗受限時能取得 2–6\% 的時間與耗電改善。最有價值的是「隨收斂門檻調整位元數」：總能量印出的每一位都不變，而且不需要事先知道安全的位數。

% ==================================================================
\section{高功耗的真實應用：DF-MP2}\label{sec:mp2}

第 \ref{sec:sci} 節的科學軟體在 1000~W 下平均用不到上限（金表面約 535–600~W、矽 1000 約 800~W），必須壓低功耗上限才看得到明顯的加速。因此我們另外選擇一個高功耗的應用：GPU4PySCF 的密度擬合二階微擾理論（DF-MP2）。它的 MP2 階段 GPU 功耗約 800~W（1000~W 時模組約 870~W，只有部分時間撞上限），矩陣乘法佔 GPU 時間的 78\%，而且幾乎都是同一種收縮：把密度擬合張量的兩個切片相乘（FP64 DGEMM）。截斷實作於程式內：張量建好後原地截斷一次尾數，之後所有矩陣乘法都受惠，不需要複製。同一個程序先做一次 HF，再依序執行各種 MP2 版本（紫杉醇，def2-TZVP 基底，每個上限交錯執行兩輪，截斷變體的兩輪差距都小於 0.6 個百分點），結果列於表~\ref{tab:mp2}。

\begin{table}[tbp]
\centering
\caption{DF-MP2（紫杉醇，def2-TZVP）各變體的 MP2 時間與相對原生的加速}\label{tab:mp2}
\small
\begin{tabularx}{\textwidth}{@{}>{\raggedright\arraybackslash}p{2.2cm}>{\centering\arraybackslash}p{1.2cm}CCCC@{}}
\toprule
 & 原生 & 截斷 32 位 & 截斷 20 位 & 內建 FP32 & cuBLAS 模擬 32 位 \\
\midrule
能量誤差 & — & $10^{-11}$\Eh & $10^{-8}$\Eh & $(3\text{–}4)\times10^{-8}$\Eh & \ee{2}{-8}\Eh \\
\addlinespace
1000~W & 39.5~s & +0.3\% & +0.8\% & +2.6\%、+11.5\%$^{*}$ & 慢 36\% \\
700~W & 44.1~s & +5.3\% & +7.6\% & +18\% & 慢 28\% \\
500~W & 52.0~s & +4.9\% & +8.3\% & +15\% & 慢 18\% \\
\addlinespace
1000~W 時少用的功耗（耗電） & — & 61–70~W（8\%） & 106–115~W（14\%） & — & — \\
\bottomrule
\end{tabularx}
\begin{tablenote}原生欄為 MP2 階段的時間，其餘為相對原生的加速；最後一列為 1000~W 時比原生少用的模組功耗與每次計算少耗的電。$^{*}$兩輪的結果，差距大。\end{tablenote}
\end{table}

截斷在功耗受限時快 4.9–8.3\%；1000~W 時速度幾乎不變（+0.3–0.8\%），但少用 61–115~W，與 HPL 同一量級。GPU4PySCF 內建的 FP32 更快，誤差也在化學精度內，因為這個張量的指數跨度很窄（整體 27 個、列內中位數 18 個 2 的次方），FP32 放得下。cuBLAS 模擬在這個形狀一律較慢。

\paragraph{小結} 在高功耗的真實應用上，截斷取得與 HPL 同一量級的收益；但資料跨度窄時，改變格式（FP32）可能更快。改變格式在什麼情況下會出問題，見第 \ref{sec:span} 節。

% ==================================================================
\section{跨度寬且功耗高的應用：高斯過程地統計}\label{sec:gp}

前面的應用不是功耗高而資料跨度窄（DF-MP2），就是資料跨度寬而功耗不高：高斯基底 DFT 的列內跨度很寬（第 \ref{sec:span} 節），但幾種 DFT 變體在 1000~W 下只用到 650–690~W。不過只要功耗上限低於它們的用電，寬跨度的 DFT 也會受限：纈胺黴素（格點等級 6）在 550／450~W 時，截斷 36 位快 3.3–3.5\%、省電 2–3\%，誤差 $10^{-6}$\Eh；模擬要用到 40 位才不會算錯，時間是截斷的 1.5 倍。

\subsection{對數概似評估}\label{sec:gp-eval}

為了找到兩者兼具的例子，我們轉向地統計的高斯過程最大概似估計（ExaGeoStat 這類應用）。每次評估對數概似都要對 $n\times n$ 的 Matérn 共變異矩陣做 Cholesky 分解，計算量與 HPL 一樣幾乎全是矩陣乘法；共變異隨距離指數衰減，所以同一列中數值大小相差很多。我們撰寫了一個 ExaGeoStat 式的核心（$n = 80000$ 個隨機位置、分塊 Cholesky、尾端更新為 FP64 矩陣乘法），在 6 個範圍參數上評估對數概似；截斷實作於程式內（每個面板在尾端更新前，以依記憶體順序走訪的單次融合核心截斷一次；早先兩次分開、存取不合併的版本只快 5.2–5.7\%）。原生在 1000~W 用到 953–970~W，面板有 88\% 的元素比所在列最大值小 $2^{-32}$ 以上，跨度寬與功耗高兩者兼具。結果列於表~\ref{tab:gp}。

\begin{table}[tbp]
\centering
\caption[高斯過程對數概似評估的加速與偏差]{高斯過程對數概似評估相對原生 FP64 的加速與偏差（$n = 80000$，6 個範圍參數，兩輪）}\label{tab:gp}
\small
\begin{tabularx}{\textwidth}{@{}l>{\raggedright\arraybackslash}p{3.5cm}>{\raggedright\arraybackslash}p{3.5cm}LL@{}}
\toprule
上限 & 截斷 36 位（偏差 $\ee{2}{-5}$） & 截斷 28 位（$\ee{7}{-4}$） & cuBLAS 模擬 36 位（$\ee{2}{-3}$） & FP32（87） \\
\midrule
1000~W & +2.4\%／+3.2\%，省電 2.6–4.5\% & +3.0\%／+3.9\%，省電 7.2–8.7\% & 慢 41\% & 慢 8.9\%／8.2\% \\
700~W & +4.9\%／+5.1\% & +8.2\%／+7.6\% & 慢 31\% & +3.5\%／+3.0\% \\
500~W & +5.0\%／+4.7\% & +7.7\%／+7.4\% & 慢 21\% & +4.3\%／+3.6\% \\
\bottomrule
\end{tabularx}
\begin{tablenote}偏差為對數概似的最大絕對偏差（對數概似的數值約為 $10^{5}$）；「／」分隔兩輪的結果。\end{tablenote}
\end{table}

截斷 28 位在功耗受限時快 7.4–8.2\%；1000~W 時這組也快 3.0–3.9\% 並省電 7–9\%，但 1000~W 的加速隨設定而變（第 \ref{sec:gp-mp} 節依 Morton 排序的 $n = 60000$ 只快 0.0–0.1\%，因為原生只有部分時間撞上限）；偏差對參數估計沒有影響。FP32 的偏差大到無法用於估計（$n = 8000$ 時 FP32 仍估出相同的值，$n$ 變大後才偏離，可能與條件數隨 $n$ 變差有關）；cuBLAS 模擬比截斷 28 位不準，又慢 21–41\%（尾端更新的 $k$ 等於區塊大小 1024，依第 \ref{sec:crossover} 節的分界，這個形狀本來就不利於模擬）。另外，這裡的列內跨度比 DFT 更寬，模擬卻沒有算錯，只是精度變差：列內跨度寬是模擬出錯的必要條件，而非充分條件。

\subsection{最大概似估計}\label{sec:gp-mle}

表~\ref{tab:gp} 是在固定的 6 個範圍參數上評估；為了確認估出的參數本身不變，我們也執行一維的估計：變異數與平滑度固定，以黃金分割搜尋範圍參數（取對數），每個變體評估 17 次對數概似（$n = 60000$，1000 與 700~W，各交錯兩輪；\path{gh200/results/gp_mle_summary.csv}）。截斷 36 位、28 位與 cuBLAS 模擬估出的範圍參數都與原生完全相同（$\hat\theta = 0.05007952$，印出的 8 位全同），最大對數概似只差 $\ee{4.8}{-6}$ 與 $\ee{2.3}{-4}$；FP32 的估計偏了 0.49\%，對數概似低 0.72（$n = 8000$ 時連 FP32 都估出同一個值）。700~W 時整個估計以截斷 28 位快 7.1–7.3\%、36 位快 4.8–4.9\%，耗電少 6.6–6.8\% 與 4.5–4.8\%；1000~W 時也快 1.8–2.3\% 與 1.6–2.2\%，耗電少 9.1–9.9\% 與 3.5–5.6\%；cuBLAS 模擬慢 32–41\%。

\subsection{與 tile 混合精度疊加}\label{sec:gp-mp}

前人的地統計混合精度（第 \ref{sec:rw-geostat} 節）是把位置依空間填充曲線排序後，把整個 tile 換成低精度格式。我們在自寫的核心中加入類似的變體：位置依 Morton 曲線排序；在尾端更新中，每個 tile 欄在 FP64 範圍以外的 tile（依 tile 範數規則決定，要求精度 $10^{-9}$，同文獻~\cite{Cao2023}）改用 FP16 輸入、FP32 累加的 tensor core。範圍參數 0.02 時有 52\% 的 tile 使用 FP16，0.08 時只有 3\%。六個範圍參數合計的結果列於表~\ref{tab:gp-mp}（$n = 60000$，ABBA 兩輪；\path{gh200/results/gp_mp_summary.csv}）。

\begin{table}[tbp]
\centering
\caption[截斷與 tile 混合精度疊加的效果]{截斷與 tile 混合精度疊加的效果（$n = 60000$，Morton 排序，倍數皆相對原生 FP64）}\label{tab:gp-mp}
\small
\begin{tabularx}{\textwidth}{@{}lL>{\centering\arraybackslash}p{1.8cm}>{\centering\arraybackslash}p{2.6cm}>{\centering\arraybackslash}p{1.8cm}@{}}
\toprule
上限 & 截斷 28 位 & 混合精度 & 混合精度 + 截斷 28 位 & 疊加後再快 \\
\midrule
1000~W & 同速，省電 14.7\%（這組原生只有 19–54\% 的時間撞上限） & 1.62 倍 & 1.65 倍 & 1.9\% \\
700~W & +7.4\% & 1.60 倍 & 1.71 倍 & 7.0\% \\
500~W & +7.2\% & 1.63 倍 & 1.73 倍 & 6.1\% \\
對數概似最大偏差 & \ee{1.3}{-4} & 0.043 & 0.043 & — \\
\bottomrule
\end{tabularx}
\end{table}

截斷可以疊加在改變格式的做法之上：把混合精度剩下的 FP64 tile 再截斷到 28 位，功耗受限時再快 6–7\%，偏差不變。把 FP16 運算元也截斷到 5 個尾數位元則沒有額外收益。ExaGeoStatCPP 的混合精度選項（\texttt{--computation=mixed}）只有列舉值、沒有實作，所以無法與原版直接比較。這組實驗的位置依 Morton 曲線排序，截斷 28 位在 700／500~W 快 7.2–7.4\%，與表~\ref{tab:gp}（未排序）一致。

\subsection{實際的 ExaGeoStat}\label{sec:gp-exa}

我們建置了 ExaGeoStatCPP（StarPU 與 Chameleon 的 tile Cholesky，cuBLAS 13.6，tile 大小 1920），以攔截工具截斷它的 FP64 tile GEMM，執行最大概似估計（$N = 60000$，每次 6 個 BOBYQA 評估；資料以原生版本產生一次，之後每次都讀取同一份；\path{gh200/results/exageostat_tdp.csv}）。它的 Matérn 共變異數在 CPU 上產生，每次評估約 9 秒，GPU 分解只需 2.3 秒（1000~W）。分解時間的比較（ABBA 兩輪平均）列於表~\ref{tab:exa}。

\begin{table}[tbp]
\centering
\caption{ExaGeoStatCPP 的 tile Cholesky 分解時間比較（ABBA 兩輪平均）}\label{tab:exa}
\small
\begin{tabularx}{\textwidth}{@{}lCCCC@{}}
\toprule
上限 & 截斷 36 位對複製基準 & 截斷 28 位對複製基準 & 截斷 28 位對未修改的原生 & SM 時脈（複製基準 $\rightarrow$ 28 位） \\
\midrule
1000~W & 快 3.0\% & 快 4.0\% & 慢 3.1\% & 1913 $\rightarrow$ 1975~MHz \\
600~W & 快 5.0\% & 快 8.0\% & 慢 1.3\% & 1166 $\rightarrow$ 1305~MHz \\
450~W & 快 5.6\% & 快 9.4\% & 慢 0.5\% & — \\
\bottomrule
\end{tabularx}
\end{table}

即使在 1000~W，分解也部分受功耗限制（GPU 650–690~W，常有功耗上限造成的降頻事件）。對數概似的相對變化最多為 $\ee{1.2}{-10}$（36 位）與 $\ee{9.8}{-8}$（28 位）。450~W 的兩輪差異大（截斷 28 位對複製基準為 +6.9\% 與 +11.9\%，SM 707~MHz 對 926~MHz），因為 CPU 產生共變異數時與 GPU 共用模組上限。攔截工具本身的複製要多花 7–9\%，與未修改的 ExaGeoStat 相比仍慢 0.5–3.1\%；要取得這個收益，截斷必須實作於函式庫內（例如 Chameleon 的 GEMM codelet）。

\subsection{小結}

在跨度寬且功耗高的矩陣乘法工作上，整個矩陣使用同一種格式的做法（FP32、INT8 模擬）不是結果錯誤就是較慢；在測試的三種整個矩陣的做法（FP32、cuBLAS 模擬 36 位、截斷）中，只有截斷又快又準，估出的範圍參數也不變。前人的 tile 混合精度快 1.6 倍，截斷還能疊加其上，在功耗受限時再快 6–7\%。在實際的 ExaGeoStat 上，截斷使分解比複製基準快 3–9\%，但攔截工具的複製成本將其抵銷，因此需要在函式庫內實作。本節的限制是：主要結果來自自寫的核心與合成資料，混合精度變體是簡化的重新實作。

% ==================================================================
\section{與 FP64 模擬及單精度的分界}\label{sec:crossover}

以較少位元計算 FP64 矩陣乘法還有另一條路：以 INT8 tensor core 模擬 FP64。我們在同一節點、同一功耗上限下，把誤差對齊到同一量級後比較速度。以小矩陣與 CPU 精確值對照確認：在相同位元數下，截斷與模擬的誤差屬於同一量級。本節的截斷都是往零截斷；應用程式用的最近捨入在相同位數下歸零的位元數相同、只多一次加法，誤差則小約 2.6 倍（約 1.4 位），所以同樣精度下最近捨入可以少保留約 1 位，分界點會稍微往有利於截斷的方向移動；我們沒有以最近捨入重做這組掃描。GEMMul8 使用 3.5.1 版（2026 年 9 月 29 日發布；截至本報告，之後的 3.5.2 只修正 HIP 編譯）。

\subsection{功耗上限改變勝負}\label{sec:crossover-cap}

圖~\ref{fig:square} 比較「方陣」（$m = n = 32768$、$k = 2048$ 的合成資料，以下沿用此稱呼）上的各種做法。GEMMul8 的功耗很高：1000~W 時就用到約 970~W，450~W 時時脈掉到 650–720~MHz，反而輸給截斷。cuBLAS 模擬在 1000~W 只用約 706~W；相反地，cuBLAS 模擬在方陣上隨上限降低而更有優勢（28 位時從 1000~W 的 1.07 倍升到 450~W 的 1.20 倍）。所以上限收緊對誰有利，取決於模擬本身的功耗；截斷在每個上限都比原生快，但幅度不隨上限單調變化（第 \ref{sec:gemm} 節）。在較粗的精度等級（$10^{-9}$）則是模擬較快；cuBLAS 模擬以約 8 位為一片，精度呈階梯式變化，截斷則能以 1 位為單位調整。

\begin{figure}[tbp]
\centering
\figfile{fig5_06_emu_caps}
\caption[方陣 FP64 矩陣乘法：截斷、INT8 模擬與單精度的比較]{方陣 FP64 矩陣乘法在 450、575 與 1000~W 下相對原生的速度。左：誤差對齊的截斷 36 位與 INT8 模擬（cuBLAS 32 位、GEMMul8 10 模數）；右：單精度做法（3xTF32、FP32，誤差未對齊）。}\label{fig:square}
\begin{figsource}Slurm 工作 1894；\path{gh200/results/m3_ozaki2_tf32x3_summary.txt}\end{figsource}
\end{figure}

\subsection{矩陣形狀決定勝負}\label{sec:crossover-shape}

我們攔截 PySCF 纈胺黴素計算的全部 45,574 次 GEMM，發現兩種形狀各佔 35\% 的運算量（圖~\ref{fig:shapes}）。模擬需要先把輸入拆片、最後再合併，這部分的成本與矩陣元素數（約 $mk + kn + mn$）成正比，而乘法本身是 $mnk$。在 $n = 300$ 的瘦長矩陣上，拆片成本與乘法相當：與截斷 36 位同精度的 cuBLAS 模擬（32 位）在 1000／575~W 只有原生 0.38／0.45 倍的速度，較粗的 28 位也只有 0.47／0.57 倍；在 $k = 76800$ 的長內積矩陣上，拆片成本被均攤，INT8 的高運算量便顯現出來。截斷沒有拆片成本，在同一個上限下對任何形狀的加速都差不多（例如 1000~W 時兩種形狀都快約 5–6\%）。因此最佳做法是依形狀選擇：瘦長矩陣用截斷，長內積矩陣用模擬。

\begin{figure}[tbp]
\centering
\figfile{fig5_07_shapes}
\caption[纈胺黴素的兩種真實矩陣形狀上截斷與模擬的速度]{1000~W 下，纈胺黴素 DFT 的兩種真實矩陣形狀（瘦長 $757248\times300\times2958$ 與長內積 $2958\times2958\times76800$）及作為參考的方陣上，截斷與模擬相對原生的速度。}\label{fig:shapes}
\begin{figsource}Slurm 工作 1894、1896、1897；\path{gh200/results/m3_realshape_summary.txt}\end{figsource}
\end{figure}

\subsection{分界點掃描與預測公式}\label{sec:crossover-sweep}

為了把「依形狀選擇」變成可用的規則，我們固定 $m = 8192$，$n$ 取 256、1024、4096、8192，$k$ 取 256 到 81920，在 1000、575 與 450~W 各量測兩輪，共 480 筆（圖~\ref{fig:crossover}）。

\begin{figure}[tbp]
\centering
\figfile{fig5_08_crossover}
\caption[截斷與 INT8 模擬的分界點掃描]{截斷與 INT8 模擬的分界點掃描（$m = 8192$）：模擬時間相對截斷保留 36 位的倍數。}\label{fig:crossover}
\begin{figsource}Slurm 工作 2005–2008（480 筆）；\path{gh200/results/crossover_*.csv}\end{figsource}
\end{figure}

截斷保留 36 位在全部 60 組（形狀 $\times$ 上限）都比原生快（450~W 為 3.8–11.9\%、575~W 為 3.6–5.9\%、1000~W 為 0.5–4.8\%）。模擬只在 $n$ 與 $k$ 都大時較快：cuBLAS 模擬需要 $n \ge 4096$ 且 $k \ge 16384$，勝出的形狀在 1000、575、450~W 分別有 5、5、2 種；但在模擬本身用不滿功耗的形狀上（例如 $n = 8192$、$k = 4096$，模擬對截斷的時間比從 1.17 降到 1.05），上限越低反而對模擬越有利，方向取決於模擬本身是否受功耗限制。GEMMul8 需要 $n \ge 4096$ 且 $k \ge 4096$，在那裡時間只要截斷的 0.47–0.80 倍，但精度低一個量級（$10^{-10}$ 對 $10^{-11}$）。以
\begin{equation}
  \frac{t_{\mathrm{emu}}}{t_{\mathrm{trunc}}} = A + B\left(\frac{1}{m} + \frac{1}{n}\right) + \frac{C}{k}
  \label{eq:crossover}
\end{equation}
擬合，GEMMul8 的決定係數 $R^{2}$ 達 0.98–0.99；若只用單一的有效大小 $1/(1/m + \allowbreak 1/n + \allowbreak 1/k)$ 擬合，則只有 0.86–0.89，表示 GEMMul8 的輸入拆片與輸出合併是兩種不同的成本。cuBLAS 模擬的 $R^{2}$ 只有 0.88–0.89，與單一有效大小的模型相近，仍有模型未捕捉的因素（例如不同形狀選到不同核心）。四個 $n$ 值分別在不同節點上量測：每格的比值是同節點配對，但形狀效應與節點差異因此無法分開，這是實驗設計上的缺陷，而不只是限制。此外每個上限只有約 20 個點擬合 3 個參數，$R^{2}$ 只能當作描述性的指標。

式~\eqref{eq:crossover} 只以 $m = 8192$ 的資料擬合。我們以先前量測、未參與擬合的兩種纈胺黴素真實形狀檢驗（1000~W 與 575~W、兩種模擬方法，共 8 個案例）：瘦長的 $757248\times300\times2958$ 預測截斷勝出、長內積的 $2958\times2958\times76800$ 預測模擬勝出，8 個案例的勝負全部判斷正確。倍數本身的誤差在 cuBLAS 模擬為 6–15\%、GEMMul8 為 24–37\%，後者較大是因為瘦長矩陣的 $m = 757248$ 遠超出擬合範圍。不過這 8 個案例都離分界很遠（模擬的時間不是截斷的一半到八成，就是兩倍以上），勝負本來就不難判斷，只是寬鬆的檢驗；公式對 $m$ 的依賴也尚未驗證。因此式~\eqref{eq:crossover} 目前只是描述這組量測的經驗式；要用來替函式庫選擇方法，還需要在分界附近、其他 $m$ 以及同一節點上重新量測與驗證。

\subsection{FP32 與 3xTF32 不適用的原因}\label{sec:crossover-fp32}

3xTF32（自行撰寫的實作，不是 CUTLASS 的版本）在方陣上最快，但它的誤差隨內積長度 $k$ 成長；截斷後的 FP64 仍以 FP64 累加，誤差與 $k$ 無關（表~\ref{tab:tf32}）。

\begin{table}[tbp]
\centering
\caption{單精度做法的矩陣乘法相對誤差隨內積長度 $k$ 的變化}\label{tab:tf32}
\small
\begin{tabularx}{\textwidth}{@{}>{\raggedright\arraybackslash}p{4.6cm}CCC@{}}
\toprule
內積長度 $k$ & 3xTF32 & TF32 & FP32 \\
\midrule
2048 & \ee{4.8}{-6} & \ee{2.6}{-4} & \ee{8.1}{-7} \\
16384 & \ee{3.8}{-5} & — & \ee{2.3}{-6} \\
95386（金表面的平面波數） & \ee{5.6}{-5} & \ee{2.7}{-4} & \ee{2.8}{-6} \\
375286（矽的平面波數） & \ee{2.2}{-4} & \ee{3.4}{-4} & \ee{1.9}{-6} \\
\bottomrule
\end{tabularx}
\end{table}

在 $4096\times4096$ 上，$k$ 每增加 8 倍，3xTF32 的誤差就增加 7.5–7.9 倍，與 tensor core 以 FP32 累加時朝零捨入的既有觀察一致~\cite{Fasi2021}。Quantum ESPRESSO 的 GEMM 內積長度就是平面波數（數十萬），所以 3xTF32 退化到與 TF32 同一量級，實測無法收斂（第 \ref{sec:sci-fp32} 節）。

\subsection{小結}

截斷與模擬沒有一方全面勝出，分界點由所需精度、矩陣形狀與功耗上限決定，兩者應依形狀互補使用。把所有 GEMM 換成單精度（FP32、3xTF32），在長內積上無法取代 FP64 截斷。

% ==================================================================
\section{指數跨度：只能以截斷處理的資料}\label{sec:span}

截斷保留每個元素自己的 11 位指數，只改變尾數；改變格式的做法都必須先處理指數：FP32 與 FP8 以整個張量的範圍為準，INT8 模擬（Ozaki 類）以每一列的最大值為基準切片，比最大值小 $2^{-k}$ 以上的元素在保留 $k$ 位時會被捨去。因此「數值大小相差多少」會限制整個矩陣共用縮放的換格式做法，而截斷不受影響。

\subsection{合成實驗}\label{sec:span-synth}

我們讓矩陣乘法（$8192\times8192\times16384$；此形狀在跨度小時是模擬較快）的每個元素，其指數在 $S$ 個 2 的次方內均勻分佈（圖~\ref{fig:span}）。

\begin{figure}[tbp]
\centering
\figfile{fig5_12_spread}
\caption[指數跨度對各方法逐元素誤差的影響]{指數跨度 $S$ 變大時，各方法逐元素相對誤差的中位數。}\label{fig:span}
\begin{figsource}Slurm 工作 2034–2036（1000／575~W，兩輪）；\path{gh200/results/range_sweep.csv}\end{figsource}
\end{figure}

截斷 36 位的誤差與加速（575~W 快 5.7–6.3\%、1000~W 快 4.4–4.5\%）完全不受跨度影響。固定位數的 cuBLAS 模擬精度變差 4–9 倍（相當於少了 2–3 位），而模擬以約 8 位為一片，必須多付一整片才能追上；保證 FP64 精度的動態模式所需位數從 58 升到 301 位，超過上限就改用原生 FP64，連同判斷跨度的額外工作，時間是截斷的 1.6–2.6 倍；FP32 在跨度 256 時溢位。GEMMul8 在這種均勻分佈的資料上幾乎不受影響。另外，若只是每列各自有一個比例（整體跨度寬、列內窄），模擬完全不受影響，但 FP32 仍然溢位。

\subsection{真實應用}\label{sec:span-apps}

我們在攔截工具中加入逐列跨度統計（每個矩陣乘法抽樣 256 列／行，並以合成資料驗證與理論值相符），量出各應用中「比所在列最大值小 $2^{-32}$ 以上」的元素比例，並在各應用上直接比較截斷與模擬（表~\ref{tab:span-apps}、圖~\ref{fig:rangemap}）。

\begin{table}[tbp]
\centering
\caption{各應用的列內指數跨度，以及截斷與 INT8 模擬的誤差}\label{tab:span-apps}
\small
\begin{tabularx}{\textwidth}{@{}>{\raggedright\arraybackslash}p{2.0cm}L>{\centering\arraybackslash}p{2.1cm}>{\centering\arraybackslash}p{3.0cm}>{\centering\arraybackslash}p{1.9cm}@{}}
\toprule
應用 & 列內比最大值小 $2^{-32}$ 以上的元素 & 截斷 36 位 & cuBLAS 模擬固定 36 位（5 片） & GEMMul8（10 模數） \\
\midrule
PySCF 纈胺黴素 DFT & 24–44\%（格點上的原子軌域矩陣達 87\%） & \ee{9.6}{-7}\Eh & \textbf{0.0525\Eh（約 33~kcal/mol）} & 發散 \\
QE 矽 & 0.24–0.54\% & 0~Ry & \ee{6}{-8}~Ry & 記憶體不足 \\
QE 金表面 & 0.008–0.02\% & 0~Ry & \ee{5}{-8}~Ry & \ee{3.7}{-7}~Ry \\
\bottomrule
\end{tabularx}
\begin{tablenote}元素比例依運算量加權。QE 金表面的跨度統計取自只跑 2 圈的基準輸入，能量誤差取自完整收斂的輸入。「記憶體不足」指 GPU 記憶體不足。\end{tablenote}
\end{table}

\begin{figure}[tbp]
\centering
\figfile{fig5_13_rangemap}
\caption[各應用的整體指數跨度與列內小元素比例]{各應用的整體指數跨度與列內小元素比例，以及固定位數的 INT8 模擬是否算錯。}\label{fig:rangemap}
\begin{figsource}\path{gh200/results/range_map_points.csv}（每個點的來源列於檔內；HPL 與高斯過程的點來自工作 3291、3304、3305）\end{figsource}
\end{figure}

同樣是模擬 36 位，在列內跨度窄的 QE 上完全正常；在列內跨度寬的 DFT 上，SCF 照常收斂、沒有任何警告，能量卻錯了 52\mEh，而截斷在兩者上都正常。模擬加到 40 位（6 片）就能算對（誤差 $\ee{8}{-8}$\Eh），但時間是截斷 36 位的約 1.5 倍（第 \ref{sec:gp} 節）；出錯的是使用者指定的固定位數模式。cuBLAS 的自動模式（依指數跨度決定位數，上限設為 79 位）則沒有算錯（誤差 $\ee{4}{-9}$\Eh）：89\% 的運算量需要 80 位以上而改用原生 FP64，其餘以 72–79 位模擬，加上判斷跨度的額外工作，時間是原生的 2.1 倍（1000~W）。所以寬跨度資料上的模擬並非不能用，而是不划算。反過來，跨度很窄的 HPL（依運算量加權，只有 $\ee{2.7}{-6}$ 的元素比所在列最大值小 $2^{-32}$ 以上）以固定 36 位模擬也通過殘差檢查（0.197；截斷 36 位為 0.057），但在單 GPU、$N = 81920$ 時只有原生 65\% 的速度。合成的均勻分佈低估了這個問題，因為真實的原子軌域值隨距離指數衰減，同一列中很多元素比最大值小得多。

\subsection{小結}

指數跨度會限制改變格式的做法：逐張量縮放的格式取決於整體跨度，逐列切片的 INT8 模擬取決於列內跨度。跨度寬時，整個矩陣共用縮放的格式或固定位數的模擬可能出錯，截斷則不受影響；但列內跨度寬只是模擬出錯的必要條件（第 \ref{sec:gp} 節的高斯過程就沒有出錯），以 tile 為單位改變格式也能處理這類資料。列內跨度可以用攔截工具事先量測，作為選擇方法時的風險指標。

% ==================================================================
\section{已使用 FP8 的 AI 工作負載}\label{sec:ai}

在 AI 上，FP8（3 個尾數位元）已經是 Hopper 的主流格式，我們從訓練與推論兩個方向測試尾數歸零，細節見附錄 \ref{app:ai}。Llama 2 70B LoRA 微調中，把凍結的 FP8 基底權重再截到 2 位，每一步只快約 1\%，驗證損失卻從第一次評估就高 0.008–0.010，達標反而慢 19–39\%（兩個 seed、四組配對）。推論中，以最近捨入保留 3 位的 BF16 與 FP8 精度相當，在文獻中 FP8 容易出問題的場景也打平，但 FP8 的矩陣乘法快 1.7–1.9 倍；在 FP8 之上再截斷，則只是以精度換取速度。因此在 Hopper 上，FP8 已經取得「減少尾數位元」幾乎全部的好處，尾數歸零的價值在 FP64 科學計算。

